\documentclass[11pt]{article}

\usepackage[margin=1in]{geometry}
\usepackage{amsmath}
\usepackage{booktabs}
\usepackage{enumitem}
\usepackage{graphicx}
\usepackage{siunitx}
\usepackage{hyperref}



\title{Crawling Flat Ground Speed Test}
\author{Rift / Team 09}
\date{\today}

\begin{document}


\maketitle

\section*{Test ID}
\textbf{MP-01 Crawling Flat Ground Speed Test} 

\section*{Test Objective}
Quantify the average forward speed of the robot on a flat surface under a fixed control command.

\section*{Robot Configuration}
\begin{itemize}[noitemsep]
    \item Robot configuration: Nominal resting shape
    \item Chamber pressures: 15 psi
    \item Control mode: Open-loop
    \item Electronics status: Batteries are fully charged
\end{itemize}

\section*{Materials}
\begin{itemize}[noitemsep]
    \item tape measures
    \item digital protractor
    \item chalk
    
\end{itemize}

\section*{Environment and Setup}
\begin{itemize}[noitemsep]
    \item Surface type: Concrete
    \item Surface angle: 0~\si{\degree}
    \item External load: None 
    \item Measurement method: stopwatch and tape measurer, camera
\end{itemize}

\section*{Initial Conditions}
\begin{itemize}[noitemsep]
    \item Initial position: $x_0 = 0$
    \item Initial orientation: $\theta_0 \pm 0~\si{\degree}$
    \item Initial shape state: Crawling Configuration
    \item Stabilization time before command: 10?~s
\end{itemize}

\section*{Command Sequence}
\begin{itemize}[noitemsep]
    \item Command type: \underline{\hspace{5cm}}
    \item Command values: \underline{\hspace{5cm}}
    \item Update rate: 3~Hz
    \item Command Steps: 8~s
\end{itemize}

\section*{Measured Quantities}
\begin{itemize}[noitemsep]
    \item Net displacement $\Delta x$
    \item Time to achieve net displacement $\Delta t$
    \item Average speed
\end{itemize}

\section*{Success Criteria (Per Trial)}
\begin{itemize}[noitemsep]
    \item Total Step Number = 8
    \item Orientation change $\leq$ 10 ~\si{\degree} 
    \item Tube Pressure $\geq$ 12 psi
\end{itemize}

\section*{Test Procedure}
\begin{enumerate}[noitemsep]
    \item Place robot in the nominal resting position
    \item Complete pretest checklist
    \item Send command to move into crawling configuration
    \item Place robot in the specified initial configuration
    \item Allow system to stabilize for the specified duration
    \item Apply the command sequence
    \item Measure distance from start line to Node 1
    \item Stop the robot and save all recorded data
\end{enumerate}

\section*{Number of Trials}
\begin{itemize}[noitemsep]
    \item Total trials: 2
    \item Time between trials: 60~s
\end{itemize}

\section*{Results}
Summarize the observed performance metrics.

\begin{center}
\begin{tabular}{lccc}
\toprule
Metric & Mean & Std. Dev. & Max \\
\midrule
Speed (cm/s) &  &  &  \\
Displacement (cm) &  &  &  \\
Orientation Change (deg) &  &  &  \\
\bottomrule
\end{tabular}
\end{center}

\section*{Observed Failure Modes}
Describe any observed failure behaviors.

\begin{itemize}[noitemsep]
    \item Slipping
    \item Stalling
    \item Excessive rotation
    \item Shape instability
\end{itemize}

\section*{Conclusions}
Provide a concise interpretation of the results.

\textit{Example:} The robot achieved consistent forward motion on flat ground with an average speed of \underline{\hspace{1cm}}~cm/s. Performance degraded beyond a slope of \underline{\hspace{1cm}}~\si{\degree}, where sustained forward progress was not achieved.

\section*{Notes and Limitations}
Document sources of variability or limitations.

\end{document}
